\section{Short solutions of the 18 exercises}
\label{sec:solutions}
$\lambda(\text{Cu K}\alpha) = \SI{1.5406}{\angstrom}$ unless stated. \readoff{} values are read from the figures ($\pm0.1$--$0.2^\circ$,
$\pm0.3$ eV, $\pm10$\,\% in size).

\newcommand{\ex}[1]{\paragraph{Ex.~#1}}

\ex{1 (UV-vis, Au--Ag)} All particles are $\approx$20 nm, so the shift is due to composition. i $\approx$400 nm = Ag;
ii $\approx$415 = Au$_{11}$Ag$_{89}$; iii $\approx$430 = Au$_{28}$Ag$_{72}$; iv $\approx$520 = Au\readoff. Class note: ``the more Ag the alloy
has, the more energetic the peak'' \slide{166}.

\ex{2 (XRD, fcc)} Peaks 28.5, 33.1, 47.5, 56.3, 59.1, 69.4, 76.7, 79.1\textdegree\readoff{} give $N$ = 3, 4, 8, 11, 12, 16, 19, 20
$\to$ (111), (200), (220), (311), (222), (400), (331), (420). $a = d\sqrt N$ = 5.41 \AA{} for every peak $\Rightarrow$
$\boxed{a\approx\SI{5.41}{\angstrom}}$ (matches CeO$_2$).

\ex{3 (BCC or FCC)} $\sin^2\theta$ ratios 1 : 2.01 : 3.02 : 4.04 : 5.05 : 6.06 $\Rightarrow$ \textbf{BCC}: (110), (200), (211), (220),
(310), (222). $a$ = 3.185, 3.178, 3.172, 3.171, 3.171, 3.169 \AA{} $\Rightarrow$ $\boxed{a\approx\SI{3.17}{\angstrom}}$ (Cu K$\alpha$
assumed; W-like\extra).

\ex{4 (Ag vs Ag$_2$O, size, XPS)} Peaks 38.1, 44.2, 64.4, 77.3\textdegree{} = fcc (111), (200), (220), (311) with
$a$ = 4.09 \AA{} $\Rightarrow$ \textbf{metallic Ag}. Ag$_2$O ($a$ = 4.723 \AA) would also show peaks at 26.7, 32.8 and 54.9\textdegree;
its (200) coincides with Ag (111) at 38.1\textdegree, so one peak is not enough.
Scherrer: $B = 0.94-0.16 = 0.78^\circ = 0.0136$ rad, $\theta = 19.05^\circ$:
$L = 0.9\cdot1.5404/(0.0136\cdot0.945)\approx\SI{108}{\angstrom}\approx\boxed{\SI{11}{nm}}$ (9 nm with quadratic subtraction\extra).
XPS: yes, from the Ag 3d chemical shift (small; Ag$_2$O is anomalously slightly \emph{lower}\extra), the lattice O 1s
($\approx$529 eV\extra) and the Ag MNN Auger peak\extra.

\ex{5 (TiO$_2$/TiN)} XPS: red, Ti 2p$_{3/2}\approx$458.6 eV = \textbf{TiO$_2$} (ref.\ 458.8 \slide{157}); blue, $\approx$455 eV =
\textbf{TiN}: N is less electronegative than O, so Ti is less positive and has lower $E_B$; the shoulder at 458 eV is surface oxide.
XRD: 36.7, 42.6, 61.8, 74.1, 78.0\textdegree\readoff{} = (111), (200), (220), (311), (222), $\boxed{a\approx\SI{4.24}{\angstrom}}$.
Nanorods: 0.21 nm fringes at 90\textdegree{} = $d_{200}$ (4.24/2 = 2.12 \AA) $\Rightarrow$ viewed along [001], exposed faces
$\boxed{\{100\}}$.

\ex{6 (CeO$_2$--ZrO$_2$)} (a) Each sample shows \textbf{one} fcc set between CeO$_2$ (28.55, 33.09, 47.49, 56.36\textdegree) and
ZrO$_2$ (30.39, 35.23, 50.68, 60.25\textdegree): A $\approx$ 29.0, 33.6, 48.5, 57.6; B $\approx$ 29.8, 34.6, 49.7, 59.1\readoff{}
$\Rightarrow$ \textbf{solid solution} Ce$_{1-x}$Zr$_x$O$_2$ ($a_A\approx5.32$, $a_B\approx\SI{5.18}{\angstrom}$; Vegard $x\approx$ 0.3 and 0.7\extra,
\checkflag{compositions depend on my readings}).
(b) FWHM(111) $\approx$2.2\textdegree{} (A) and 0.7\textdegree{} (B)\readoff{} $\Rightarrow$ $L_A\approx\SI{4}{nm}$, $L_B\approx\SI{12}{nm}$.
Stirring homogenises supersaturation $\to$ many nuclei $\to$ small crystallites; without stirring, fewer nuclei grow larger.
(c) HRTEM $d$ per particle (one intermediate $d$ vs 3.12 or 2.94 \AA); STEM-EDX/EELS maps (Ce and Zr in the same particle?).

\ex{7 (Pd/Al$_2$O$_3$)} Digital images 2\_1a, 2\_1b, 2\_2 are not in the folder: apply recipe G ($\geq$100 particles, scale bar
20/5/20 nm). Sample A: small, well-dispersed particles. Epitaxy: check whether Pd fringes (2.25/1.95 \AA) are aligned with
the alumina fringes in 2\_1b.

\ex{8 (CeO$_2$ cubes)} Images 3\_1/3\_2 are not in the folder; expected answer: at 573 K, a cube bounded by \{100\}, with
\{200\} fringes (2.71 \AA) parallel to the edges (\{111\}, 3.12 \AA, at an angle). At 973 K the particle has sintered and is rounded,
with \{111\} (and \{110\}) facets at the truncated edges. Verify by measuring $d$.

\ex{9 (Pd@NiO microreactor)} (a) 1.95 \AA{} = Pd(200) (1.946); Ni gives 2.035/1.762 \AA $\Rightarrow$ \textbf{Pd core, NiO shell}.
(b) EDX/EELS maps, XPS Ni 2p (sees the shell: Ni$^{2+}$ + satellites), HRTEM/XRD of NiO ($d_{200}$ 2.09 \AA\extra).
(c) Core 10 nm, total 50 nm\readoff: $t_\mathrm{core}$ = 10/2 = 5 min at 35\,\textdegree C, 30 \si{\micro L/min} $\to$ 150 \si{\micro L};
$t_\mathrm{shell}$ = 40/10 = 4 min at 5\,\textdegree C, 60 \si{\micro L/min} $\to$ 240 \si{\micro L}. With $r$ = 0.5 mm
($A$ = 0.785 mm$^2$): $L$ = \textbf{19 cm} and \textbf{31 cm}; add KMnO$_4$ at the end (\cref{fig:micro}).
\checkflag{the handwritten equations on slide 102 (``NaBH$_2$ + 2HCl'', ``KMnO$_3$'') are electron bookkeeping, not balanced reactions.}
\begin{figure}[h]
\centering
\includegraphics[width=0.55\linewidth]{microreactor_class.png}
\caption{Class solution of the microreactor exercise (slide 102).}
\label{fig:micro}
\end{figure}

\ex{10 (Cu/CeO$_2$)} (a) A: Cu 2p$_{3/2}\approx$934 eV with strong satellites (938--945 eV) = Cu$^{2+}$ (CuO)\extra; B: sharp 932.5 eV,
no satellites $\Rightarrow$ \textbf{Cu is reduced to metallic Cu} during reaction (Cu$^+$ would need the LMM Auger peak to exclude\extra).
(b) 3.12 \AA{} = CeO$_2$(111) (3.124) = support; 1.81 \AA{} = Cu(200) (1.808) = particle.
(c) The Cu fringes are parallel to the ceria fringes at the interface $\Rightarrow$ epitaxial (Cu\{200\}$\parallel$CeO$_2$\{111\}).
(d) Before: CuO, so no plasmon band; after: metallic Cu, so a band at $\approx$570 nm (weak and broad for small particles); ceria absorbs
below 400 nm. (e) Impregnation of CeO$_2$ with Cu(NO$_3$)$_2$ $\to$ calcination (CuO) $\to$ H$_2$ reduction.

\ex{11 (preparations)} (a) Ag triangles: capping agent adsorbing on \{111\} (or solvothermal). (b) Pt@SiO$_2$: Pt colloid + sol--gel
TEOS shell, in a microreactor (class note). (c) TiO$_2$ photonic crystal: polystyrene-sphere template + Ti precursor, then
calcine (ordered macropores, \slide{45}). (d) Carbon aerogel: organic sol--gel (resorcinol--formaldehyde\extra) $\to$
supercritical drying $\to$ pyrolysis in inert gas\extra.

\ex{12 (CoMg cubes; Pd/SiO$_2$ aerogel)} (A) capping agents on \{100\} or hydrothermal (NaOH, $T$; class note: ``either capping
agents or hydrothermal''). (B) TEOS sol--gel with Pd salt/colloid in the sol $\to$ gel $\to$ supercritical drying $\to$
H$_2$ reduction (or impregnate the aerogel afterwards; liquid may collapse pores\extra).

\ex{13 (Au@SiO$_2$)} The plasmon grows with core size: black (clear 520 nm band) = (b) largest core; red (weak band) = (a);
blue (no band) = (c) tiny core. Controlled preparation: microreactor with Turkevich core (size via citrate ratio and time),
then a TEOS sol--gel shell (thickness via TEOS concentration and time).

\ex{14 (Au sizes)} Mean sizes $\approx$7.5, 11 and 19 nm\readoff, all below 30 nm, where the plasmon barely moves (518--522 nm), and the
distributions overlap $\Rightarrow$ \textbf{hardly distinguishable}; at most a slight red shift and narrower band for the largest.
Use TEM or XRD widths instead.

\ex{15 (Au nanotubes + Pt)} (a) Galvanic replacement of Ag nanowires with HAuCl$_4$ ($E = 1.50-0.80>0$); Pt by reduction of
H$_2$PtCl$_6$ (NaBH$_4$/ascorbic acid) on the tubes. (b) 0.237 nm $\approx d_{111}$(Au) = 0.235 nm $\Rightarrow$ \{111\} walls;
metallic Au $\Rightarrow$ plasmon band (red-shifted for hollow/elongated tubes\extra). (c) TEM tomography; XPS (Pt inside is attenuated
by the Au wall); also SEM. (d) Pt$^0$ doublet at $\approx$71.4/74.7 eV, Pt$^{2+}$ at $\approx$72.6/75.9 eV\readoff. Au (2.54) is more
electronegative than Pt (2.28) $\Rightarrow$ Pt at the interface loses electron density and appears oxidised (Pt$^{\delta+}$).

\ex{16 (Au--Ag core--shell)} One conclusion per technique:
\begin{itemize}
  \item \textbf{EELS maps}: yellow fills the centre, red forms a ring $\Rightarrow$ core--shell; faceted shapes.
  \item \textbf{XRD}: fcc, $a\approx4.08$ \AA; Au and Ag differ by 0.3\,\% $\Rightarrow$ XRD \emph{cannot} distinguish them.
  \item \textbf{HRTEM}: 2.35 \AA{} = (111) of either metal; also not conclusive.
  \item \textbf{UV-vis}: two bands ($\approx$400 + $\approx$490 nm) $\Rightarrow$ segregated metals; after shell removal one band at
        $\approx$540 nm $\Rightarrow$ the \textbf{core is Au}.
  \item \textbf{XPS Au 4f}: 83.8 eV as prepared (below 84.0) $\Rightarrow$ Au gains charge from Ag (2.54 vs 1.93), so they are in contact or partly alloyed;
        84.0 eV after removal = pure Au.
  \item \textbf{XPS Ag 3d}: main component $\approx$368.6 eV ($\delta^+$ Ag bonded to Au) + minor $\approx$368.2 eV (pure Ag)\readoff.
\end{itemize}
\textbf{Answer:} \textbf{Au core @ Ag shell} with partial Au--Ag alloying at the interface, not a homogeneous alloy.
Preparation: seeded growth: Au cores with a capping agent (shapes), then reduce AgNO$_3$ on them.

\ex{17 (capped Pt)} AFM height $\approx$5 nm, TEM core $\approx$2.5 nm\readoff{} $\Rightarrow$ the capping shell is still there:
$\ell\approx(5-2.5)/2\approx\boxed{\SI{1.2}{nm}}$. AFM widths are useless (tip convolution).

\ex{18 (Rh/TiO$_2$)} A (small, flat) and B (medium): Rh fringes aligned with TiO$_2$ $\Rightarrow$ \textbf{epitaxy}; C (large, round,
small contact) has its own orientation $\Rightarrow$ no epitaxy. For catalysis choose \textbf{A}: highest dispersion and longest
metal--support perimeter, and anchored against sintering.
